% Section 2.6, from lectures/L02/appendix/e_exercises.md.

\section{Exercises}
\label{c2:sec:exercises}

\begin{aside}
\textbf{How to check your work.} Exercise~\ref{c2:ex:driver} is checked by this chapter's test
suite: write the file at the path the specification gives, then run \code{make test}. The suite is
cumulative, so it runs Chapter~1's \code{shift_bits} tests too; if something you write here
breaks it, this is where it shows up.

The rest will have worked solutions in the companion repository, published later as
\file{lectures/L02/appendix/f_solutions.md}, and the hand calculations have short answers in
Appendix~\ref{app:answers}. Work each one before you read it.

\textbf{No hardware is needed for any of this.} \Secref{c2:sec:board} is optional and nothing
depends on it.
\end{aside}

Each exercise is labelled with its kind. There is exactly one \textbf{Cross-check}, and here it is
Exercise~\ref{c2:ex:crosscheck}.

% ------------------------------------------------------------------------------------------
\exercise{Four states}{Recall}
\label{c2:ex:states}
\begin{parts}
  \item Give the four combinations of \code{DDRx} and \code{PORTx} for a single pin, and say what
    each one does to that pin.
  \item Which of the four is a mistake rather than a choice? Say what a program reading it gets.
  \item \code{PINx} is not in that table. Give both of the things it does, and say why neither of
    them is what its name suggests.
  \item You want to toggle one output pin. Give the one-instruction way and the three-instruction
    way, and one reason to prefer the short one that is not about speed.
\end{parts}

% ------------------------------------------------------------------------------------------
\exercise{Wiring a button}{Design}
\label{c2:ex:button}
A button is wired between an ATmega328P pin and ground, with nothing else attached.
\begin{parts}
  \item Give the values of \code{DDRx} and \code{PORTx} for that pin, and say what each one is
    doing.
  \item What does \code{PINx} read when the button is not pressed? When it is?
  \item A driver subroutine \code{btn_pressed} should return 1 when the button is down. Write, in
    words, what it has to do to the bit it read, and say why.
  \item Somebody proposes leaving \code{PORTx} at 0 and fitting an external pull-down resistor to
    ground, wiring the button to 5\,V instead. Their polarity is now the intuitive one. Give one
    real advantage of their scheme and two reasons nobody does it.
  \item Your driver reads the pin every millisecond and counts a press each time it sees a 1 then
    a 0. Presses are being counted several times each. Explain, without using the word
    ``software''.
\end{parts}

% ------------------------------------------------------------------------------------------
\exercise{Translating pin numbers}{Hand calculation}
\label{c2:ex:translate}
For Arduino pins \textbf{3}, \textbf{8} and \textbf{13}, give the port, the bit number, the mask,
and the data space addresses of the three port registers. Do this from \secref{c2:sec:pinmap} and
\secref{c1:sec:iowindow}, not from a table you have already been shown.
\begin{parts}
  \item Fill in the three rows.
  \item What is the mask for pin 8, and what would it be if you forgot to subtract 8 first? Say
    what the LED does in each case.
  \item Arduino pins stop at 13. Port B has eight bits. What are the other two, and why can you not
    use them?
\end{parts}

\noindent\textbf{Check yourself:} the port and bit for every pin are in \secref{c2:sec:pinmap}, and
deriving a row is quicker than looking it up.

% ------------------------------------------------------------------------------------------
\exercise{What a call leaves behind}{Recall}
\label{c2:ex:call}
\begin{parts}
  \item An \code{rcall} at byte address \code{0x0A} calls a subroutine. What two things does the
    instruction do, in order?
  \item Exactly what value goes on the stack? Say whether it is a byte address or a word address,
    and give the number.
  \item Which byte of it ends up at the higher address?
  \item \code{SP} was \code{0x08FF}. What is it inside the subroutine, and what is it after
    \code{ret}?
  \item After the \code{ret}, what is stored at \code{0x08FF}?
  \item A subroutine does one \code{push} and no \code{pop}, then \code{ret}. Describe what
    happens, and say why there is no error message.
\end{parts}

% ------------------------------------------------------------------------------------------
\exercise{The LED driver}{Code}
\label{c2:ex:driver}
Write \file{drivers/source/led.asm} to the specification in \secref{c2:sec:driver}.
\begin{parts}
  \item Write all five subroutines. \code{.include "led.inc"} for the field offsets.
  \item Run \code{make test}.
  \item Disassemble your \code{led_init} and count how many instructions run in the case where the
    pin number is rejected. Compare it against the case where it is accepted.
  \item Deliberately move the range check to \emph{after} the structure is filled in, rebuild, and
    say which test catches it. Then put it back.
\end{parts}

% ------------------------------------------------------------------------------------------
\exercise{Read, modify, write}{Design}
\label{c2:ex:rmw}
\begin{parts}
  \item Write the three-instruction sequence that sets one bit of a port register, given the mask
    in \code{r24} and the register's address in \code{X}.
  \item Somebody replaces it with a single \code{st X, r24}. Their LED works. Describe precisely
    what else on that port stops working, and say why their own test did not catch it.
  \item Which test in this chapter's suite does catch it? Look it up.
  \item \code{led_toggle} needs no read-modify-write at all. Say why, and give the one property
    that gives it besides being shorter.
  \item The AVR has \code{sbi} and \code{cbi}, which set and clear one bit of a low I/O register in
    a single instruction. Give the reason this driver cannot use them, and note that it has
    nothing to do with speed.
\end{parts}

% ------------------------------------------------------------------------------------------
\exercise{Predicting the driver's cost}{Hand calculation}
\label{c2:ex:predict}
You know from Chapter~1 that \code{shift_bits} costs $6n + 10$ cycles and
\code{shift_bits_inverted} costs $7n + 10$, where $n$ is the number of shifts.
\begin{parts}
  \item Without counting the rest of \code{led_on}, say how much more it must cost for Arduino pin
    13 than for pin 8, and show why.
  \item Do the same for \code{led_off}.
  \item \code{led_init} calls both shift routines once each. How much more does it cost for pin 13
    than for pin 8?
  \item Which of the five subroutines is the cheapest for a given pin, and why?
  \item Explain why you can answer all of the above without knowing how the rest of each
    subroutine is written.
\end{parts}

% ------------------------------------------------------------------------------------------
\exercise[fillaccent]{A driver whose cost depends on its argument}{Cross-check}
\label{c2:ex:crosscheck}
\emph{Compute it by hand, measure it, reconcile.}

Do the parts in order and write each answer down before starting the next.
\begin{parts}
  \item \textbf{By hand.} From Exercise~\ref{c2:ex:predict}, predict the difference in cycles
    between \code{led_on} for Arduino pin 13 and \code{led_on} for pin 8. Convert it to
    microseconds at 16\,MHz.
  \item \textbf{In full, on paper.} Write the count out in the shape \secref{c1:sec:onpaper} gives,
    once for bit 0 and once for bit 5. Your table has to include the nested call, so it needs a row
    for the \code{rcall} and then every row of \code{shift_bits} itself. Do this before (c).
  \item \textbf{By measurement.}
\begin{shell}
make measure CALLS="led_init:0x0200:8 led_on:0x0200"
make measure CALLS="led_init:0x0200:13 led_on:0x0200"
\end{shell}
    Two calls, because \code{led_on} needs a structure that \code{led_init} had to build first.
  \item \textbf{Reconcile.} The \textbf{difference} between the two pins should come out the same
    in both. If your absolute figures disagree while the difference agrees, say what that tells you
    about where the error is, and find it: a constant offset is a row you left out of the table,
    and it is the same row for both pins.
  \item \textbf{Now the asymmetry.} Predict whether \code{led_enabled} takes the same number of
    cycles when the LED is on as when it is off. Write the prediction down, then measure both:
\begin{shell}
make measure CALLS="led_init:0x0200:13 led_on:0x0200 led_enabled:0x0200"
make measure CALLS="led_init:0x0200:13 led_off:0x0200 led_enabled:0x0200"
\end{shell}
  \item \textbf{The discrepancy.} State the difference in cycles and explain it in one sentence.
    Then say what it would take to write a version of \code{led_enabled} whose cost did
    \textbf{not} depend on the answer, and whether you think it is worth doing.
  \item \textbf{What it means.} You are bit-banging a protocol that needs a pulse held for exactly
    5\,µs, and you are using this driver. Say which of the five subroutines you could use inside
    that pulse and which you could not, and give the one change to the \emph{wiring} that would
    make the problem smaller.
\end{parts}

% ------------------------------------------------------------------------------------------
\exercise{A program that blinks}{Code}
\label{c2:ex:blink}
Extend \file{drivers/app/main.asm} so that it initialises the LED on pin 13 and toggles it
forever.
\begin{parts}
  \item Write it. There is no delay routine yet, so it will toggle as fast as the loop runs.
  \item Work out how fast that is, in cycles per toggle, from your \code{led_toggle} cost and the
    loop's own \code{rjmp}. Convert it to a frequency.
  \item An LED toggling at that frequency does not look like it is blinking. Say what it looks like
    and why, and what would have to change for it to blink visibly. Chapter~5 is where you fix it
    properly.
  \item \code{make build} should report that it built \file{app.hex}. If you own a board,
    \secref{c2:sec:board} will put it on one.
\end{parts}
